\subsection{Three additive weights characterize compressing discs exactly}
\label{sec:three-weight-criterion}

Let $C$ be a connected component of the supplied proper normal surface.
The cell-incidence weight gives $\chi(C)$. Let $\gamma_1,\gamma_2$ be the
checked homology basis of the boundary torus described above, and define
\[
 h_i(C)=\sum_{e\in\gamma_i}|C\cap e|\pmod2,
 \qquad i=1,2.
\]
The raw orbit calculation retains the corresponding integer sums; parity is
taken only when interpreting the resulting component signature. In
particular, signed Euler arithmetic and binary homology arithmetic are not
silently identified.

\begin{lemma}
\label{lem:boundary-character-completeness}
The pair $(h_1(C),h_2(C))$ vanishes exactly when the mod-two boundary class of
$C$ vanishes in the ambient torus.
\end{lemma}
\begin{proof}
On every boundary triangle, a normal arc contributes two endpoints.
Consequently the parity of edge intersections is a cellular $1$-cocycle,
including in generalized triangulations where repeated face edges must be
combined modulo two. This cocycle represents the Poincar\'e dual of the
boundary multicurve. Evaluation on a homology basis determines its cohomology
class, since evaluation gives the nondegenerate pairing between
$H^1(T^2;\F_2)$ and $H_1(T^2;\F_2)$. Both spaces have dimension two. Thus
the two values vanish precisely for the zero boundary class.
\end{proof}

\begin{theorem}[Exact three-weight disc criterion]
\label{thm:three-weight-disc}
A connected component $C$ is a compressing disc for the ambient boundary
torus if and only if
\[
 \chi(C)=1,
 \qquad (h_1(C),h_2(C))\ne(0,0).
\]
\end{theorem}
\begin{proof}
A compressing disc has Euler characteristic one. Its boundary is an embedded
essential circle on the torus. Such a circle represents a primitive integral
slope, so its reduction modulo two is nonzero. The lemma proves necessity.

Conversely, the nonzero boundary class implies that $C$ has boundary. The
classification of connected compact surfaces now leaves only a disc. Indeed,
an orientable surface of genus $a$ with $b\ge1$ boundary circles has
$\chi=2-2a-b$; equality to one forces $a=0,b=1$. A nonorientable surface with
$k\ge1$ crosscaps and $b\ge1$ boundary circles has
$\chi=2-k-b\le0$. The resulting disc has nonzero boundary class, and hence
its simple boundary circle is essential on the torus.
\end{proof}

The boundary condition in this argument is necessary. A closed projective
plane also has Euler characteristic one; its two boundary evaluations are
zero, so it is correctly excluded. Such closed one-sided surfaces occur in
the supported general torus-boundary manifold fixtures, even though those
fixtures need not be knot exteriors. No assumption of irreducibility is
smuggled into the local component test.

Two binary linear probes are also the minimum needed to recognize the zero
element of the full torus boundary-class space by linear evaluation alone.
A single linear functional on the two-dimensional space has a nonzero
kernel. The two-probe statement concerns this unrestricted homology interface;
it does not prove a lower bound for a specially restricted collection of
normal surfaces in one fixed manifold, nor for nonlinear preprocessing.
Together with the Euler coordinate, it explains the natural three-weight
implementation and its reduction from $7t$-dimensional coordinate extraction.

The torus hypothesis cannot be dropped from the equivalence. On a boundary
surface of genus at least two, an essential separating simple circle can have
zero mod-two homology. The same predicate would then give only a sufficient
test for certain compressing discs, and could miss others. The native input
validator's one-torus-boundary contract is therefore part of the theorem,
not merely a convenience for constructing a small basis.

Finally, the test must be applied to individual component weights. Two
parallel meridians have zero total boundary class but contribute two
compressing discs. Conversely, an inessential vertex-link disc together with
a M\"obius band can have total Euler characteristic one and nonzero total
boundary class while contributing no compressing disc. A weighted orbit
histogram resolves both ambiguities without listing the components one by one.
